% =============================================================================
% 09_APPENDIX — Technical Appendix: AI-Assisted Literature Pipeline
% =============================================================================

\section*{Technical Appendix: AI-Assisted Literature Pipeline}
\label{sec:technical_appendix}
\addcontentsline{toc}{section}{Technical Appendix}

% =============================================================================
% A. Pipeline Architecture
% =============================================================================

\subsection{Pipeline Architecture}
\label{sec:app_pipeline}

The literature collection pipeline operates across five functional layers:
\textbf{data collection}, \textbf{paper scoring}, \textbf{LLM-based
summarization}, \textbf{storage and indexing}, and \textbf{dissemination}.
Figure~\ref{fig:pipeline_overview} illustrates the end-to-end flow.

\begin{figure}[htbp]
\centering
% Simple ASCII-flowchart-style diagram
\begin{tabular}{cl}
\textbf{Step 1} & \texttt{collector.py} queries arXiv, PubMed, Semantic Scholar, DBLP \\
\textbf{Step 2} & Results deduplicated by DOI / title and passed to \texttt{ranker.py} \\
\textbf{Step 3} & \texttt{ranker.py} computes composite priority scores \\
\textbf{Step 4} & Priority papers passed to \texttt{summarizer.py} (LLM) \\
\textbf{Step 5} & Structured summaries stored in JSON + SQLite \\
\textbf{Step 6} & Daily digest + weekly synthesis written to \texttt{outputs/} \\
\textbf{Step 7} & Telegram bot delivers digest to registered channel \\
\end{tabular}
\caption{Literature pipeline execution steps.}
\label{fig:pipeline_overview}
\end{figure}

The pipeline is implemented as two entry-point scripts:
\texttt{daily\_pipeline.py} (or \texttt{daily\_pipeline\_v2.py}) runs on a
scheduled cron job each weekday morning; \texttt{weekly\_synthesis.py} runs
every Friday. Both scripts support a \texttt{--push} flag to commit outputs
directly to the GitHub repository.

% =============================================================================
% B. Data Sources and Queries
% =============================================================================

\subsection{Data Sources and Query Strategy}
\label{sec:app_sources}

Four bibliographic sources are queried on each pipeline run:

\begin{itemize}[itemsep=0.2em]
\item \textbf{arXiv} (\url{http://export.arxiv.org/api/query}) — preprints from
  cs.AI, cs.LG, cs.CL, cs.CY, q-bio, eess.AS, cs.CV. Eight structured queries
  cover each of the five thematic axes using \texttt{all:} field operators.
  Results parsed from Atom XML response.

\item \textbf{PubMed} (\url{https://eutils.ncbi.nlm.nih.gov/entrez/eutils}) — MEDLINE
  via E-search / E-fetch. Eight PubMed-formatted queries covering clinical AI
  evaluation, governance, adaptive regulation, evidence synthesis, and
  participatory engagement.

\item \textbf{Semantic Scholar} (\url{https://api.semanticscholar.org/graph/v1})
  — free-tier Graph API. Five keyword queries with health-article filtering.
  Fields requested: paperId, title, authors, abstract, year, venue, externalIds.

\item \textbf{DBLP} (\url{https://api.dblp.org/v1}) — targeted venue searches for
  FAccT, AIES, CHI, ICLR, and NeurIPS proceedings filtered to health-AI-relevant
  titles.
\end{itemize}

Query strings are defined in \texttt{collector.py} as constants
\texttt{ARXIV\_QUERIES}, \texttt{PUBMED\_QUERIES},
\texttt{SEMANTIC\_SCHOLAR\_QUERIES}, and \texttt{DBLP\_VENS}. Deduplication
is performed by paper ID and normalized title (80-character window, alphanumeric
stripped) against the existing corpus before insertion.

% =============================================================================
% C. Paper Scoring Algorithm
% =============================================================================

\subsection{Paper Scoring Algorithm}
\label{sec:app_scoring}

Each collected paper is scored by \texttt{ranker.py} using a five-component
composite formula. The total score determines the priority band assigned to
the paper, which governs whether it is included in the daily digest and whether
an LLM summary is generated.

\subsubsection{Composite Score Formula}

{\centering
Score = Theme\_Score + Novelty\_Score + Recency\_Boost + Venue\_Score + Methods\_Bonus
\par}

Each component is computed as follows:

\begin{itemize}
\item \textbf{Theme\_Score} — Each of the five thematic keyword groups
  contributes a weight when its keywords appear in the paper title or abstract:
  \texttt{AI\_EVAL\_VALID} (weight 2), \texttt{PARTICIPATORY\_GOV} (weight 2),
  \texttt{ADAPTIVE\_REGULATION} (weight 2), \texttt{EVIDENCE\_IMPLEMENT}
  (weight 2), \texttt{CLINICAL\_HEALTH\_AI} (weight 1). Pre-tagged themes from
  the collector are also counted. Theme\_Score = matched\_themes\_count $\times$ 2.

\item \textbf{Novelty\_Score} — Detected via linguistic markers in title and
  abstract. New contribution indicators (\textit{we propose, novel, first to,
  new framework}) yield 2.0. Extension indicators (\textit{extend, adapt,
  improve}) yield 1.0. Redundant or incremental work yields 0.5.

\item \textbf{Recency\_Boost} — Based on days since publication date:
  $\leq$30 days = 2.0; $\leq$90 days = 1.0; $\leq$365 days = 0.5.

\item \textbf{Venue\_Score} — Top-tier venues (FAccT, AIES, NeurIPS, ICML,
  ICLR, Nature Medicine, JAMA, Lancet, Science, Medical Internet Research,
  Implementation Science) yield 3.0. Recognized venues (Frontiers, BMJ Global
  Health, Nature Biomedical Engineering, JAMA Network Open, IEEE J Biomed
  Health Inform) yield 1.5. Preprints (arXiv, medRxiv, bioRxiv) yield 0.5.

\item \textbf{Methods\_Bonus} — Empirical/lab studies (RCT, prospective cohort,
  validation study) yield 2.0. Systematic reviews and meta-analyses yield 1.5.
  Frameworks and conceptual proposals yield 1.0. Commentary and editorials
  yield 0.5.
\end{itemize}

\subsubsection{Priority Bands}

\begin{threeparttable}
\caption{Priority band thresholds and pipeline behaviour.}
\label{tab:priority_bands}
\begin{tabularx}{\linewidth}{ccl}
\toprule
\textbf{Priority} & \textbf{Score range} & \textbf{Pipeline behaviour} \\
\midrule
3 (Must Read) & $\geq$ 7.0 & Daily digest highlight; LLM summary generated \\
2 (Important) & 4.0 -- 6.9 & Daily digest listing; no LLM summary \\
1 (Supplementary) & 1.5 -- 3.9 & Database only \\
0 (Low) & $<$ 1.5 & Database only; excluded from outputs \\
\bottomrule
\end{tabularx}
\end{threeparttable}

The full scoring implementation is in \texttt{ranker.py}, functions
\texttt{score\_paper()} and \texttt{rank\_papers()}.

% =============================================================================
% D. Database Schema
% =============================================================================

\subsection{Database Schema}
\label{sec:app_db_schema}

Papers are persisted in two formats: a JSON document store
(\texttt{data/papers.json}) and a SQLite database (\texttt{data/papers.db}).
The JSON store holds the canonical scored corpus used by digest generation.
The SQLite database (defined in \texttt{database.py}) supports structured
queries and is initialised with the following schema:

\begin{verbatim}
CREATE TABLE papers (
    id                INTEGER PRIMARY KEY AUTOINCREMENT,
    paper_id          TEXT UNIQUE NOT NULL,
    source            TEXT NOT NULL,
    title             TEXT NOT NULL,
    authors           TEXT,
    abstract          TEXT,
    url               TEXT,
    pdf_url           TEXT,
    published_date    DATE,
    venue             TEXT,
    doi               TEXT,
    keywords          TEXT,          -- JSON array
    score             REAL DEFAULT 0.0,
    priority          INTEGER DEFAULT 0,
    summary           TEXT,
    critique          TEXT,
    methods           TEXT,
    gaps              TEXT,
    related_papers    TEXT,          -- JSON array
    processing_status TEXT DEFAULT 'pending',
    created_at        TIMESTAMP DEFAULT CURRENT_TIMESTAMP,
    last_updated      TIMESTAMP DEFAULT CURRENT_TIMESTAMP
);

CREATE TABLE daily_digests (
    id                INTEGER PRIMARY KEY AUTOINCREMENT,
    digest_date       DATE NOT NULL UNIQUE,
    papers_included   TEXT,          -- JSON array of paper_ids
    total_papers_count INTEGER,
    highlight_papers  TEXT,          -- JSON array of top papers
    trends_observed   TEXT,
    gaps_identified   TEXT,
    generated_at      TIMESTAMP DEFAULT CURRENT_TIMESTAMP,
    sent_to_telegram  BOOLEAN DEFAULT 0,
    sent_to_email     BOOLEAN DEFAULT 0,
    sent_to_gdocs     BOOLEAN DEFAULT 0
);
\end{verbatim}

Indexes are created on \texttt{score}, \texttt{published\_date}, and
\texttt{processing\_status}.

% =============================================================================
% E. Output Products
% =============================================================================

\subsection{Output Products}
\label{sec:app_outputs}

Three classes of output are produced by the pipeline:

\begin{itemize}
\item \textbf{Daily digests} (\texttt{outputs/digests/digest\_YYYY-MM-DD.md}) —
  Markdown summaries of all papers collected in a given day, organized by
  priority. Includes title, authors, source, abstract excerpt, thematic tags,
  and composite score. If a digest already exists for the same date, the
  pipeline creates a versioned file (\texttt{digest\_YYYY-MM-DD\_v2.md}, etc.)
  rather than overwriting.

\item \textbf{Weekly synthesis reports} (\texttt{outputs/weekly/weekly\_YYYY-MM-DD.md})
  — Aggregates and cross-references papers collected over the preceding seven
  days. Includes a re-ranked list of the top papers, a thematic distribution
  table, and a narrative synthesis of emerging trends.

\item \textbf{Topic-trend matrices} (\texttt{outputs/topic\_trends/matrix\_YYYY-MM-DD.md})
  — A 4-week rolling window matrix showing the proportion of papers assigned
  to each of the five thematic axes over time, enabling identification of
  shifts in the literature's thematic focus.
\end{itemize}

All outputs are committed to the GitHub repository
(\texttt{Arkoys/ai-health-lit-review}) using a Fine-Grained Personal Access
Token with read-write access to that repository.

% =============================================================================
% F. Telegram Bot Integration
% =============================================================================

\subsection{Telegram Bot Integration}
\label{sec:app_telegram}

Dissemination to the research team is automated via a Telegram bot. The bot
is configured through two environment variables (loaded from
\texttt{~/.hermes/.env} at runtime, \emph{not} the project's \texttt{.env}):

\begin{verbatim}
TELEGRAM_BOT_TOKEN=<bot_token>
TELEGRAM_HOME_CHANNEL=<home_channel_id>
\end{verbatim}

The bot delivers the daily digest summary to the configured Telegram channel
on each successful pipeline run. Delivery is triggered from
\texttt{reporter.py} and handles its own retry logic. The Hermes gateway's
Telegram integration (defined in \texttt{~/.hermes/config.yaml}) manages the
underlying Bot API calls.

% =============================================================================
% G. Cron Job Configuration
% =============================================================================

\subsection{Cron Job Configuration}
\label{sec:app_cron}

The pipeline is executed on a scheduled basis via two cron jobs managed
through the Hermes cron system (\texttt{mcp\_cronjob}):

\begin{threeparttable}
\caption{Scheduled literature pipeline cron jobs.}
\label{tab:cron_jobs}
\begin{tabular}{lllp{6cm}}
\toprule
\textbf{Job ID} & \textbf{Name} & \textbf{Schedule} & \textbf{Script} \\
\midrule
d7d9439fa180 & daily-ai-health-literature-digest & \texttt{0 9 * * 1-5} &
  \texttt{python3 daily\_pipeline\_v2.py --push} \\[4pt]
d1d4915d28d0 & weekly-ai-health-literature-synthesis & \texttt{0 11 * * 5} &
  \texttt{python3 weekly\_synthesis.py --push} \\
\bottomrule
\end{tabular}
\end{threeparttable}

The daily job runs at 09:00 Monday--Friday (weekdays only). The weekly
synthesis runs at 11:00 every Friday. Both jobs are configured with
\texttt{deliver: origin}, routing output back to the originating Telegram DM.
The \texttt{--push} flag ensures outputs are committed to GitHub before the job
completes.

Cron jobs can be manually triggered with:
\begin{verbatim}
mcp_cronjob(action='run', job_id='d7d9439fa180')
mcp_cronjob(action='run', job_id='d1d4915d28d0')
\end{verbatim}

% =============================================================================
% H. Key Implementation Scripts
% =============================================================================

\subsection{Key Implementation Scripts}
\label{sec:app_scripts}

\begin{itemize}
\item \texttt{collector.py} — Multi-source paper fetcher (arXiv, PubMed,
  Semantic Scholar, DBLP). Exports \texttt{collect\_all(dry\_run=False)} and
  \texttt{dedup\_papers(papers, existing\_ids, existing\_titles)}.

\item \texttt{ranker.py} — Scoring engine. Exports \texttt{score\_paper(paper)},
  \texttt{rank\_papers(papers)}, \texttt{score\_daily\_papers(papers, date)},
  and \texttt{generate\_ranking\_report(db, days\_back=30)}.

\item \texttt{summarizer.py} — LLM-powered paper summarizer. Generates structured
  summaries (title, authors, year, source, AI system, evaluation method,
  governance dimension, main finding, thematic tags, priority score, critique)
  for Priority 3 papers using Gemini (or OpenRouter fallback).

\item \texttt{database.py} — SQLite wrapper. Exports \texttt{PaperDatabase}
  class with \texttt{add\_paper()}, \texttt{get\_papers\_needing\_summarization()},
  \texttt{update\_paper\_summary()}, and \texttt{get\_statistics()}.

\item \texttt{daily\_pipeline.py} / \texttt{daily\_pipeline\_v2.py} — Main daily
  entry points. \texttt{daily\_pipeline\_v2.py} uses the extended four-source
  collector. Both support \texttt{--push} to commit outputs to GitHub.

\item \texttt{weekly\_synthesis.py} — Weekly aggregation. Reads digests from
  \texttt{outputs/digests/}, re-ranks papers from the database, writes
  \texttt{outputs/weekly/weekly\_YYYY-MM-DD.md} and
  \texttt{outputs/topic\_trends/matrix\_YYYY-MM-DD.md}.

\item \texttt{reporter.py} — Output delivery. Formats and sends digest summaries
  to Telegram (primary) and email (secondary).
\end{itemize}

% =============================================================================
% I. File Structure
% =============================================================================

\subsection{Project File Structure}
\label{sec:app_filestructure}

\begin{verbatim}
ai-health-lit-review/
|-- sections/
|   |-- 00_abstract.tex
|   |-- 01_introduction.tex
|   |-- 02_related_work.tex
|   |-- 03_methods.tex          <-- main text: 3-sentence pipeline summary
|   |-- 04_results.tex
|   |-- 04_thematic.tex
|   |-- 05_discussion.tex
|   |-- 05_evaluation.tex
|   |-- 06_conclusion.tex
|   |-- 07_acknowledgements.tex
|   |-- 08_appendix.tex         <-- Main Appendix (PRISMA, bib, keywords)
|   |-- 09_appendix.tex         <-- THIS FILE: Technical Appendix
|-- collector.py                <-- Multi-source paper collector
|-- ranker.py                   <-- Composite scoring algorithm
|-- summarizer.py               <-- LLM summarization
|-- database.py                 <-- SQLite wrapper
|-- daily_pipeline.py           <-- Daily pipeline entry point
|-- daily_pipeline_v2.py        <-- Extended 4-source daily pipeline
|-- weekly_synthesis.py         <-- Weekly synthesis
|-- reporter.py                 <-- Telegram / email delivery
|-- config.yaml                 <-- Sources, keywords, venues, providers
|-- .env                        <-- GITHUB_TOKEN, GITHUB_REPO
|-- data/
|   |-- papers.json             <-- Canonical scored corpus (gitignored)
|   |-- papers.db               <-- SQLite store (gitignored)
|-- outputs/
|   |-- digests/                <-- Daily digests (gitignored)
|   |-- weekly/                 <-- Weekly syntheses (gitignored)
|   |-- topic_trends/           <-- Rolling topic matrices (gitignored)
|   |-- reviews/                <-- Ad hoc paper review notes
|-- references.bib
|-- report.tex
|-- requirements.txt
|-- venv/                       <-- Virtual environment (currently symlinked)
\end{verbatim}

Note: The \texttt{data/} and \texttt{outputs/} directories are gitignored.
Outputs are pushed explicitly via \texttt{git add -f <file>} using the
Fine-Grained PAT with repository write access.

% =============================================================================
% J. Configuration
% =============================================================================

\subsection{Configuration Reference}
\label{sec:app_config}

LLM provider selection is configured in \texttt{config.yaml} under the
\texttt{providers} key. The primary provider is OpenRouter
(\texttt{stepfun/step-3.5-flash:free}) with an OpenRouter fallback
(\texttt{nvidia/nemotron-3-super-120b-a12b:free}). Additional providers
(Ollama, OpenAI, Gemini, HuggingFace) are available but disabled by default.

The five thematic keyword groups and priority venue lists are defined in both
\texttt{config.yaml} and \texttt{ranker.py}. These must be kept in sync when
updating search terms or venue tiers.
